\section{Method}

Many group-based RLVR methods \cite{shao2024deepseekmath, liu2025understandingr1zeroliketrainingcritical, xie2025unlocking, le2025promptleftbehindexploiting} primarily focus on the design of advantage estimators, while paying comparatively little attention to the more fundamental problem of reward estimation that underlies advantage computation. In practice, computational and memory constraints limit the number of sampled responses per prompt, making accurate reward estimation inherently difficult.

In contrast, we introduce a new perspective that reformulates RLVR through the lens of statistical estimation (\cref{subsec:problem_formulation}). Specifically, we view the reward not as a deterministic signal, but as a distribution induced by the policy, and frame advantage computation as a problem of estimating this distribution from finite samples. Building on this reformulation, we propose \textbf{D}iscounted \textbf{B}eta--\textbf{B}ernoulli (\textbf{DBB}) reward estimation for RLVR, which estimates the non-stationary reward distribution by leveraging historical rewards (\cref{subsec:grpohistbeta}). Despite introducing bias, DBB reduces estimator variance and achieves lower MSE. Moreover, it fundamentally prevents variance collapse and preserves informative training signals (\cref{subsec:comparison}).


\subsection{Reward Estimation as Distributional Inference}
\label{subsec:problem_formulation}

For a given prompt $q \sim \mathcal{D}$ and training step $\eta$ with corresponding epoch $\tau$, each response generated by the policy $\pi_{\theta_{\eta,\mathrm{old}}}$ produces a binary reward indicating whether the response is correct. Under our formulation, this reward is naturally modeled as a Bernoulli random variable 
\begin{equation}
X_{\tau,i} \sim \mathrm{Bernoulli}\!\left(p_\tau\right),
\label{eq:bern}
\end{equation}
where 
$p_\tau = \mathbb{P}\!\left(X_{\tau,i}=1 \mid q, \pi_{\theta_{\eta,\mathrm{old}}}\right)$
denotes the probability of obtaining a correct response.

Within this framework, group-based RLVR algorithms such as GRPO and Dr.GRPO can be interpreted as implicitly estimating the reward distribution from a finite set of rollout samples. Given $N$ observed rewards $\{X_{\tau,i}\}_{i=1}^{N}$, existing approaches rely on \textit{point estimation} of $p_\tau$ via the empirical mean 
\begin{equation}
\hat{p}_\tau^\mathrm{pt} = \frac{1}{N} \sum_{i=1}^{N} X_{\tau,i}.
\end{equation}

The corresponding reward variance is then estimated using the sample variance,
\begin{equation}
\begin{aligned}
\widehat{\mathrm{Var}}^{\mathrm{pt}}(X_{\tau})
=\;
\frac{\sum_{i=1}^{N}
\left(X_{\tau,i} - \hat{p}_\tau^{\mathrm{pt}}\right)^2}{N-1}
=
\frac{N \cdot \hat{p}_\tau^{\mathrm{pt}}
\bigl(1 - \hat{p}_\tau^{\mathrm{pt}}\bigr)}{N-1}\,
.
\end{aligned}
\end{equation}

Since the estimated variance $\widehat{\mathrm{Var}}^{\mathrm{pt}}$ is determined by the point estimator (mean) $\hat{p}_\tau^\mathrm{pt}$, it suffices to consider only the point estimator in our analysis.
Concretely, the point estimator $\hat{p}_\tau^\mathrm{pt}$ has expectation and variance as follows:
\begin{equation}
\mathbb{E}[\hat{p}_\tau^\mathrm{pt} \mid p_\tau] = p_\tau,
\end{equation}
\begin{equation}
\mathrm{Var}(\hat{p}_\tau^\mathrm{pt} \mid p_\tau) = \frac{p_\tau(1-p_\tau)}{N}.
\label{eq:var_pt}
\end{equation}
While the point estimator is unbiased, its variance grows as the number of rollouts $N$ decreases. As a result, individual rollout outcomes can exert disproportionate influence on the estimate, causing the estimator to fluctuate widely. Moreover, the estimated variance can collapse to zero when all sampled rewards are identical. Since group-relative methods compute advantages directly from these reward estimates, such estimation noise naturally propagates to the advantage signal, leading to unstable policy updates.


\subsection{Discounted Beta--Bernoulli Reward Estimation}
\label{subsec:grpohistbeta}

To estimate the Bernoulli reward distribution beyond empirical averaging, we adopt a well-established Bayesian perspective that leverages the conjugacy between the Beta prior and the Bernoulli likelihood. In this setting, the reward probability is modeled using a Beta prior and updated via a Beta posterior after observing rollout outcomes, leading to the naive Beta--Bernoulli model defined below.

\begin{definition}[Beta--Bernoulli Reward Model]
\label{def:beta_bernoulli}
To model uncertainty in the Bernoulli reward distribution, we place a Beta prior over the reward probability:
\begin{equation}
p_\tau \sim \mathrm{Beta}(\alpha_\tau, \beta_\tau),
\end{equation}
where $(\alpha_\tau, \beta_\tau)$ denote the posterior parameters for a prompt at epoch $\tau$, with initialization $\alpha_0=\beta_0=1$. 

Given $N$ rollouts with $S_\tau$ successes, the posterior distribution is given by:
\begin{equation}
\label{eq:param_update_bb}
\alpha_{\tau} = \alpha_{\tau-1} + S_\tau,
\qquad
\beta_{\tau} = \beta_{\tau-1} + N - S_\tau.
\end{equation}
\end{definition}

However, the above model assumes a stationary reward distribution. In RLVR, this assumption does not hold: due to the on-policy nature of training, the policy $\pi_{\theta}$ evolves over time, inducing a non-stationary reward distribution. As a result, historical observations may become outdated and should not be weighted equally with recent rollouts. 


\begin{algorithm}[t]
\caption{Discounted Beta--Bernoulli Reward Estimation}
\label{alg:grpohistbeta}
\begin{algorithmic}[1]
\INPUT LLM policy $\pi_\theta$, training dataset $\mathcal{D}$
\STATE Initialize $(\alpha^q_0,\beta^q_0) \leftarrow (1,1) \quad \forall q\in\mathcal{D}$

\FOR{$\text{epoch} \ \tau = 1,\dots,I$}
  \FOR{$\text{iteration} = 1,\dots,M$}
    \STATE Sample a minibatch $\mathcal{D}_b \subset \mathcal{D}$
    \STATE Set old policy $\pi_{\theta_{\mathrm{old}}} \leftarrow \pi_{\theta}$
    \FOR{each prompt $q \in \mathcal{D}_b$}
      \STATE Sample outputs $\{o_i\}_{i=1}^N \overset{\text{i.i.d.}}{\sim} \pi_{\theta_{\mathrm{old}}}(\cdot \mid q)$
      \STATE Compute rewards $\{X_i=r(o_i, q)\}_{i=1}^N$
      \STATE Update posterior parameters:
      \STATE \hspace{1em}\(
      \begin{aligned}
        \alpha^q_{\tau} &\leftarrow \lambda \alpha^q_{\tau-1} + \sum_{i=1}^N X_i, \\
        \beta^q_{\tau}  &\leftarrow \lambda \beta^q_{\tau-1}  + N - \sum_{i=1}^N X_i
      \end{aligned}
      \)
      \STATE Compute advantage $\hat{A}_i$ with $(\alpha^q_{\tau},\beta^q_{\tau})$ and $X_i$
    \ENDFOR

    \FOR{$\text{update} = 1,\dots,U$}
      \STATE Update $\pi_{\theta}$ by maximizing the objective (Eq.~\ref{eq:grpo})
    \ENDFOR
  \ENDFOR
\ENDFOR
\end{algorithmic}
\end{algorithm}


To address this challenge, we introduce the \textbf{D}iscounted \textbf{B}eta--\textbf{B}ernoulli (\textbf{DBB}) reward model, which gradually discounts past information to estimate the reward distribution accurately. 

\begin{definition}[Discounted Beta--Bernoulli Reward Model]
\label{def:discounted_beta_bernoulli}
Given posterior parameters $(\alpha_{\tau-1}, \beta_{\tau-1})$ and $N$ rollouts with $S_\tau$ successes at epoch $\tau$, the update is defined as
\begin{equation}
\label{eq:param_update_dbb}
\alpha_{\tau} = \lambda \alpha_{\tau-1} + S_\tau,
\qquad
\beta_{\tau} = \lambda \beta_{\tau-1} + N - S_\tau,
\end{equation}
where the discount factor $\lambda \in (0,1]$ controls the influence of historical observations. 
\end{definition}

Under the DBB reward model, we estimate the mean and variance of the Bernoulli reward distribution as follows:
\begin{equation}
\hat{p}_\tau^\mathrm{dbb}
=
\frac{\alpha_\tau}{\alpha_\tau + \beta_\tau},
\end{equation}
\begin{equation}
\widehat{\mathrm{Var}}^{\mathrm{dbb}}(X_{\tau} \mid \alpha_{\tau}, \beta_{\tau})
=
\frac{\alpha_\tau\beta_\tau}{(\alpha_\tau + \beta_\tau)^2}.
\end{equation}

To understand the statistical behavior of the DBB estimator, we analyze its expectation and variance:
\begin{equation}
\mathbb{E}\!\left[\hat{p}_\tau^{\mathrm{dbb}} \mid p_\tau, \alpha_{\tau-1}, \beta_{\tau-1}\right] =
w\, \mu_{\tau-1} + (1-w)\, p_\tau,
\end{equation}
\begin{equation}
\mathrm{Var}\!\left(\hat{p}_\tau^{\mathrm{dbb}} \mid p_\tau, \alpha_{\tau-1}, \beta_{\tau-1} \right)
=
(1-w)^2\,\frac{p_\tau(1-p_\tau)}{N},
\end{equation}
where $\mu_{\tau-1} = \frac{\alpha_{\tau-1}}{\alpha_{\tau-1}+\beta_{\tau-1}}$ denotes the historical posterior mean and 
$w = \frac{\lambda(\alpha_{\tau-1}+\beta_{\tau-1})}{\lambda(\alpha_{\tau-1}+\beta_{\tau-1})+N}$ controls the contribution of past observations.

These show that the DBB estimator introduces bias through shrinkage toward the historical mean, but substantially reduces variance relative to the point estimator (\cref{eq:var_pt}). Importantly, unlike estimated variance based on small-sample rollouts, the posterior variance in DBB cannot collapse to zero, ensuring stable and informative reward signals throughout training. Algorithm~\ref{alg:grpohistbeta} summarizes the overall training procedure of the group-based RLVR algorithm with DBB.

% \begin{table}[htbp]
\centering
\renewcommand{\arraystretch}{2.0}   
\setlength{\tabcolsep}{4pt}         

\caption{Comparison between point estimation and DBB estimation for reward distribution modeling. We denote the Bernoulli distribution by $\mathrm{Bern}(\cdot)$.}
\label{tab:estimator_comparison}
\begin{tabular}{l|c|c}
\toprule
 & {\small\textbf{Point Estimation}} & {\small\textbf{DBB Estimation}} \\
\midrule

$X_{\tau,i}$
& $X_{\tau,i} \sim \mathrm{Bern}(p_\tau)$
& $X_{\tau,i} \sim \mathrm{Bern}(p_\tau)$ \\


$p_\tau$
& fixed but unknown
& $p_\tau \sim \mathrm{Beta}(\alpha_\tau,\beta_\tau)$ \\

$\hat{p}_\tau$
& $\displaystyle \hat{p}_\tau^{\mathrm{pt}} = \frac{1}{N}\sum_{i=1}^N X_{\tau,i}$
& $\displaystyle \hat{p}_\tau^{\mathrm{dbb}} = \frac{\alpha_\tau}{\alpha_\tau+\beta_\tau}$ \\

$\widehat{\mathrm{Var}}(X_\tau)$
& $\displaystyle \frac{N}{N-1}\,\hat{p}_\tau^{\mathrm{pt}}\!\left(1-\hat{p}_\tau^{\mathrm{pt}}\right)$
& $\displaystyle \hat{p}_\tau^{\mathrm{dbb}}\!\left(1-\hat{p}_\tau^{\mathrm{dbb}}\right)$ \\


$\mathbb{E}[\hat{p}_\tau]$
& $p_\tau$
& $w\mu_{\tau-1} + (1-w)p_\tau$ \\

$\mathrm{Bias}(\hat{p}_\tau)$
& $0$
& $w(\mu_{\tau-1} - p_\tau)$ \\

$\mathrm{Var}(\hat{p}_\tau)$
& $\displaystyle \frac{p_\tau(1-p_\tau)}{N}$
& $\displaystyle (1-w)^2\frac{p_\tau(1-p_\tau)}{N}$ \\

\bottomrule
\end{tabular}
\end{table}

% Algorithm~\ref{alg:grpohistbeta} summarizes the overall training procedure of the group-based RLVR algorithm with DBB, and Table~\ref{tab:estimator_comparison} summarizes and compares the statistics of point estimation and ours.


\subsection{Mean Squared Error of the DBB estimator}
\label{subsec:comparison}

% 실제로 우리 DBB estimator 가 reward distribution 예측에 얼마나 효과적인지를 분석하기 위해 MSE 구해봄. underlying reward probability 에 어떻게 영향을 받는지 확인하기 위해서 앞선 분석들에서 alpha, beta 로 표현된 식을 p1~p_tau 로 구성된 수식으로 재전개해볼 것.

To assess the effectiveness of our DBB estimator in predicting the reward distribution, we compute its mean squared error (MSE). To make the estimator's dependence on the underlying reward probabilities $\{p_1, \ldots, p_\tau\}$, we re-express its expectation and variance, originally formulated in terms of the posterior parameters $(\alpha_{\tau-1}, \beta_{\tau-1})$.

Based on the update defined in \cref{eq:param_update_dbb}, posterior parameters can be expressed as
\[
\alpha_\tau = \lambda^\tau \alpha_0 + \sum_{k=1}^{\tau} \lambda^{\tau-k} S_k,
\;
\beta_\tau = \lambda^\tau \beta_0 + \sum_{k=1}^{\tau} \lambda^{\tau-k} (N - S_k).
\]
For convenience, define the sum of the posterior parameters as $H_\tau$:
\begin{equation}
H_\tau = \alpha_\tau + \beta_\tau = \lambda^\tau(\alpha_0+\beta_0) + N\sum_{k=1}^{\tau}\lambda^{\tau-k}.
\end{equation}

Under our formulation (\cref{eq:bern}), the number of successful rollouts at epoch $k$ satisfies $S_k \sim \mathrm{Binomial}(N, p_k)$. Given this, the expectation and variance of the DBB estimator can be written as follows:
\begin{equation}
\mathbb{E}\!\left[\hat{p}_\tau^{\mathrm{dbb}} \mid p_{1:\tau} \right]
= \sum_{k=0}^{\tau} c_k p_k,
\label{eq:hb_expectation_pseq}
\end{equation}
\begin{equation}
\mathrm{Var}\!\left(\hat{p}_\tau^{\mathrm{dbb}} \mid p_{1:\tau} \right)
=
\frac{\sum_{k=1}^{\tau} \lambda^{2(\tau-k)} N p_k(1-p_k)}{H_\tau^2},
\label{eq:hb_variance_pseq}
\end{equation}
where $p_0=\frac{\alpha_0}{\alpha_0+\beta_0}$ is the reward probability of the initial prior, and the weights $c_0=\frac{\lambda^\tau(\alpha_0 + \beta_0)}{H_\tau}$ and $c_k=\frac{N \lambda^{\tau-k}}{H_\tau}$ are constants. A detailed derivation is provided in Appendix~\ref{app:hb-derivation}.

Combining Equations~\eqref{eq:hb_expectation_pseq} and~\eqref{eq:hb_variance_pseq}, the MSE of the DBB estimator under non-stationary rewards is
\begin{equation}
\mathrm{MSE}\!\left(\hat{p}_\tau^{\mathrm{dbb}} \mid p_{1:\tau} \right)
=
\left(\mathbb{E}\!\left[\hat{p}_\tau^{\mathrm{dbb}} \mid p_{1:\tau} \right] -p_\tau \right)^2
+
\mathrm{Var}\!\left(\hat{p}_\tau^{\mathrm{dbb}} \mid p_{1:\tau} \right).
\label{eq:hb_mse_pseq}
\end{equation}

In contrast, the MSE of the point estimator depends only on the current reward distribution
\begin{equation}
\mathrm{MSE}\!\left(\hat{p}_\tau^{\mathrm{pt}} \mid p_\tau \right)
=
\frac{p_\tau(1-p_\tau)}{N},
\label{eq:pt_mse_pseq}
\end{equation}
and is therefore more prone to high variance, especially in low-sample settings.

% Equation~\eqref{eq:hb_mse_pseq} shows that the benefit of the DBB estimation depends on how close the historical reward probabilities $\{p_k\}_{k<\tau}$ are to the current $p_\tau$.
% When reward distribution shifts are moderate relative to the variance of the point estimator, the variance reduction obtained by using historical evidence outweighs the induced bias, yielding a strictly lower MSE.
% We demonstrate this empirically in Section~\ref{sec:mse_analysis} and provide a detailed derivation for Equations~\eqref{eq:hb_expectation_pseq} and~\eqref{eq:hb_variance_pseq} in Appendix~\ref{app:hb-derivation}.

The bias of the DBB estimator becomes smaller when the underlying reward probabilities evolve gradually, as past values $\{p_k\}_{k<\tau}$ remain close to the current $p_\tau$. 
In contrast, the variance depends on both the discount factor $\lambda$ and the number of rollouts $N$. 
On the other hand, the point estimator relies solely on $p_\tau$ and is therefore more susceptible to high variance, particularly in low-sample settings.
Considering these factors, we empirically evaluate the resulting MSE and present the findings in Section~\ref{sec:mse_analysis}.


% So far, the analysis of the DBB estimator has been conditioned on the historical posterior parameters $(\alpha_{\tau-1}, \beta_{\tau-1})$.
% We now reformulate the mean squared error (MSE) of the estimator directly in terms of the underlying reward probabilities $\{p_1, \ldots, p_\tau\}$, which makes the dependence on reward distribution shift explicit.

% Recall that at each epoch $k$, rollout rewards satisfy
% $S_k \sim \mathrm{Binomial}(N, p_k)$, Equation~\eqref{eq:param_update_dbb} can be reformulated as
% \[
% \alpha_\tau = \lambda^\tau \alpha_0 + \sum_{k=1}^{\tau} \lambda^{\tau-k} S_k,
% \;
% \beta_\tau = \lambda^\tau \beta_0 + \sum_{k=1}^{\tau} \lambda^{\tau-k} (N - S_k).
% \]

% The DBB mean estimator is
% \[
% \hat{p}_\tau^{\mathrm{dbb}} = \frac{\alpha_\tau}{H_\tau},
% \]
% where the total mass $H_\tau \triangleq \alpha_\tau + \beta_\tau = \lambda^\tau(\alpha_0+\beta_0) + N\sum_{k=1}^{\tau}\lambda^{\tau-k}$, which is deterministic given $\lambda$ and $N$.

% Conditioned on the reward probabilities $p_{1:\tau}$, its expectation can be written as
% \begin{equation}
% \mathbb{E}\!\left[\hat{p}_\tau^{\mathrm{dbb}} \mid p_{1:\tau}\right]
% =
% \rho_0 \mu_0 + \sum_{k=1}^{\tau} \rho_k p_k,
% \label{eq:hb_expectation_pseq}
% \end{equation}
% where $\rho_0 \triangleq \frac{\lambda^0(\alpha_0+\beta_0)}{H_0}$,
% $\rho_k \triangleq \frac{N\lambda^{\tau-k}}{H_\tau}$, 
% and $\mu_0 \triangleq \frac{\alpha_0}{\alpha_0+\beta_0}$, all of which are deterministic.

% Thus, the estimator mean is an exponentially weighted average of past reward probabilities.

% The corresponding conditional variance is given by
% \begin{equation}
% \mathrm{Var}\!\left(\hat{p}_\tau^{\mathrm{dbb}} \mid p_{1:\tau}\right)
% =
% \frac{\sum_{k=1}^{\tau} \lambda^{2(\tau-k)} N p_k(1-p_k)}{H_\tau^2}.
% \label{eq:hb_variance_pseq}
% \end{equation}

% Combining Equations~\eqref{eq:hb_expectation_pseq} and~\eqref{eq:hb_variance_pseq}, the conditional MSE of the DBB estimator under a non-stationary reward process is
% \begin{equation}
% \mathrm{MSE}\!\left(\hat{p}_\tau^{\mathrm{dbb}} \mid p_{1:\tau}\right)
% =
% \left(\mathbb{E}\!\left[\hat{p}_\tau^{\mathrm{dbb}} \mid p_{1:\tau}\right]\right)^2
% +
% \mathrm{Var}\!\left(\hat{p}_\tau^{\mathrm{dbb}} \mid p_{1:\tau}\right).
% \label{eq:hb_mse_pseq}
% \end{equation}

% In contrast, the point estimator depends only on the current reward distribution and has
% \begin{equation}
% \mathrm{MSE}\!\left(\hat{p}_\tau^{\mathrm{pt}} \mid p_\tau\right)
% =
% \frac{p_\tau(1-p_\tau)}{N}.    
% \label{eq:pt_mse_pseq}
% \end{equation}

% Equation~\eqref{eq:hb_mse_pseq} shows that the benefit of DBB estimation depends on how close the historical reward probabilities $\{p_k\}_{k<\tau}$ are to the current $p_\tau$.
% When reward distribution shifts are moderate relative to the variance of the point estimator, the variance reduction obtained by using historical evidence outweighs the induced bias, yielding a strictly lower MSE.
% We demonstrate this empirically in Section~\ref{sec:mse_analysis} and provide a detailed derivation for Equations~\eqref{eq:hb_expectation_pseq} and~\eqref{eq:hb_variance_pseq} in Appendix~\ref{app:hb-derivation}.
